\begin{claim}{D.2.40}{Proved}{The fourth-loop five-type and the failure of pure descent}
Put $B=K(\FF_2,2)$, $Y=\Sigma^2B$, $Z'=Y\langle7\rangle$,
$E=\Omega^2Z'$ and $U=\Omega^2E$.  Thus $Z'$ is six-connected
and $U$ is two-connected.  Let $\rho_2$ denote reduction modulo two
and $i:\FF_2\to\mathbb Z/4$ the inclusion $i(1)=2$.
There is a space-level equivalence
\[
 P_5U\simeq
 \operatorname{fib}\bigl(i\operatorname{Sq}^2\rho_2:
      K(\mathbb Z/4,3)\longrightarrow K(\mathbb Z/4,5)\bigr)
       \times K(\FF_2,5).
 \tag{D.2.40a}
\]
This statement does not assert an infinite-loop splitting.

Use the free double-loop homology coordinates of D.2.29 and write
\[
 \mu_{11}=\sigma^{*2}\bigl((\xi_1k)^*+(\xi_1f)^*\bigr)
       \in H^{11}(U;\FF_2).
 \tag{D.2.40b}
\]
Here the two ambient functionals on $H_{13}(\Omega^2Y;\FF_2)$
are restricted to $E$ before the two cohomology suspensions, and
$f\in H_6B$ is dual to $u^3$ in Serre's polynomial presentation.
Then
\[
 \mu_{11}\notin
 \operatorname{im}\bigl(H^{11}(P_5U;\FF_2)
                       \longrightarrow H^{11}(U;\FF_2)\bigr).
 \tag{D.2.40c}
\]
This excludes pure descent to the five-type.  It does not determine
the restriction of $\mu_{11}$ to $V=U\langle6\rangle$, or its
evaluation on the finite relative source in D.2.38.
\end{claim}
\begin{proof}
All unlabeled cohomology has coefficients $\FF_2$.
We first determine the two Postnikov invariants and then compute
only the degree-eleven quotient needed for \textup{(D.2.40c)}.

\medskip\noindent
\textbf{The two Postnikov invariants.}
The actual target calculations of D.2.18--D.2.22 give
\[
 \pi_3U=\mathbb Z/4\{\alpha\},\qquad
 \pi_4U=\mathbb Z/4\{\rho\},\qquad
 \alpha\eta_3=2\rho\ne0,\qquad
 \pi_5U=\mathbb Z/2\{\lambda\}.
 \tag{D.2.40d}
\]
Serre's integral calculation gives
\[
 H_4(K(\mathbb Z/4,3);\mathbb Z)=0,
 \qquad H_5(K(\mathbb Z/4,3);\mathbb Z)=\mathbb Z/2.
 \tag{D.2.40e}
\]
Indeed, his stable groups $A_n$ equal $H_{q+n}$ for $q>n$;
for a cyclic group of order at least four and $n>0$ they equal
those for $\mathbb Z/2$, indexed by admissible sequences of
degree $n$ whose first entry is even.  For $n=1$ there is none,
and for $n=2$ only $(2)$; take $q=3$.
See \cite[\S2, nos.~14--15, pp.~207--209,
Theorems~8--9]{I_Serre1953}.
The coefficient universal-coefficient theorem therefore identifies
\[
 H^5(K(\mathbb Z/4,3);\mathbb Z/4)
       =\operatorname{Hom}(\mathbb Z/2,\mathbb Z/4)
       \cong\mathbb Z/2.
 \tag{D.2.40f}
\]
Its nonzero class is $i\operatorname{Sq}^2\rho_2$:
$\operatorname{Sq}^2\rho_2$ is the unique nonzero mod-two
degree-five class, and coefficient inclusion is nonzero on
the displayed $\operatorname{Hom}$ group.
The zero Postnikov invariant would give a product of the two
Eilenberg--Mac Lane spaces and hence $\alpha\eta_3=0$.
Thus \textup{(D.2.40d)} forces the unique nonzero invariant and
identifies the first factor $P=P_4U$ in \textup{(D.2.40a)}.

The actual mod-two Hurewicz image of $\lambda$ is nonzero.
Two homology suspensions carry it to the nonzero class
$R=\{a,c\}\in H_7E$, as identified in D.2.19 and the
$D_0\vee S^7_R$ model of D.2.22.
In the next Postnikov fibration
\[
 K(\FF_2,5)\longrightarrow P_5U\longrightarrow P_4U,
 \tag{D.2.40g}
\]
the base is simply connected and the fiber map is an isomorphism
on $\pi_5$, since $\pi_5P_4U=\pi_6P_4U=0$.
A nonzero invariant in $H^6(P_4U;\FF_2)$ would evaluate
nontrivially on a mod-two six-cycle.  Its homology transgression
would kill the unique fiber $H_5$ class, forcing the Hurewicz
image of that $\pi_5$ generator to be zero in the total space.
This contradicts the value just established.  The invariant is
therefore zero, proving the space-level splitting.

\medskip\noindent
\textbf{The low Serre calculation.}
Write $x_3$ for the mod-two fundamental class of
$K(\mathbb Z/4,3)$, $y_4$ for its higher coefficient connecting
class, and $z_5=\operatorname{Sq}^2x_3$.
The coefficient fibration also gives the exact alternative model
\[
 P\simeq\operatorname{fib}\bigl(
 \operatorname{Sq}^2\rho_2+\operatorname{Sq}^1:
 K(\mathbb Z/4,3)\times K(\FF_2,4)
           \longrightarrow K(\FF_2,5)\bigr).
 \tag{D.2.40h}
\]
In fact $K(\mathbb Z/4,4)\to K(\FF_2,4)
\xrightarrow{\operatorname{Sq}^1}K(\FF_2,5)$ is the coefficient
fibration, classified by $i$; pulling it back gives precisely the
fiber of $i\operatorname{Sq}^2\rho_2$.
Consequently there is an actual class $r_4$ on $P$, restricting
to the fiber's mod-two fundamental class, with
\[
 \operatorname{Sq}^1r_4=z_5.
 \tag{D.2.40i}
\]
This is an actual Bockstein extension, not an inference from
the dimensions of a spectral sequence.

For the required spanning assertion use the Serre sequence of
$K(\mathbb Z/4,4)\to P\to K(\mathbb Z/4,3)$.
Serre's polynomial-generator theorem gives the complete fiber
generator list through degree eleven:
\[
\begin{gathered}
 r_4,\ s_5,\ \operatorname{Sq}^2r_4,\
 \operatorname{Sq}^3r_4,\ \operatorname{Sq}^2s_5,\
 \operatorname{Sq}^3s_5,\ \operatorname{Sq}^4s_5,\\
 \operatorname{Sq}^4\operatorname{Sq}^2r_4,\
 \operatorname{Sq}^5\operatorname{Sq}^2r_4,\
 \operatorname{Sq}^4\operatorname{Sq}^2s_5.
\end{gathered}
 \tag{D.2.40j}
\]
Here $s_5$ is the higher coefficient connecting class for
$\mathbb Z/4$, not $\operatorname{Sq}^1r_4$ in the total space.
The simple multiplicative system also includes $r_4^2,s_5^2$
in degrees eight and ten.
Serre's universal path-fibration construction makes these
generators and their $2^j$ powers transgressive, with suspension
and Steenrod operations compatible with transgression;
see \cite[\S2, no.~7, Theorem~1 and paragraph~7.3,
pp.~201--203; no.~11, Theorem~4 and its proof,
pp.~204--205]{I_Serre1953}.
After pullback a universally transgressive class survives every
earlier differential by naturality.  If its transgression is
zero, it is permanent, since no later outgoing differential can
start on that fiber row.

The degree-five fundamental class of $K(\mathbb Z/4,5)$ pulls
back to zero, while its higher degree-six coefficient class
pulls back to $x_3^2$.  For the latter assertion, pulling
$0\to\FF_2\to\mathbb Z/8\to\mathbb Z/4\to0$ back along
$i$ gives $0\to\FF_2\to\mathbb Z/4\to\FF_2\to0$;
thus the coefficient identity is $\delta_4i=\operatorname{Sq}^1$.
It follows that the complete nonzero differential list in the
needed range is
\[
 d_6(s_5)=x_3^2,\qquad
 d_{10}(\operatorname{Sq}^4s_5)=z_5^2.
 \tag{D.2.40k}
\]
Indeed Cartan and $\operatorname{Sq}^1x_3=0$ give
\[
 \operatorname{Sq}^2(x_3^2)=\operatorname{Sq}^3(x_3^2)
 =\operatorname{Sq}^5(x_3^2)=0,\quad
 \operatorname{Sq}^4(x_3^2)=z_5^2,\quad
 \operatorname{Sq}^4\operatorname{Sq}^2(x_3^2)=0.
 \tag{D.2.40l}
\]
In particular $s_5^2$, a separate transgressive simple-system
generator, survives.  The two nonzero images are a regular
sequence in the relevant polynomial algebra: $x_3^2$ is a
non-zero-divisor, and after quotienting by it the independent
polynomial generator $z_5$ still has non-zero-divisor square.
Thus these differentials remove the odd powers of $s_5$ and
$\operatorname{Sq}^4s_5$, retain $s_5^2$, and quotient by
$(x_3^2,z_5^2)$.  There are no additional annihilator cycles.
Every other generator needed through total degree eleven is
permanent by the universal-transgression argument, and
multiplicativity controls all mixed terms in that range.

The complete base generator list through degree eleven is
\[
\begin{gathered}
 x_3,\ y_4,\ z_5,\ \operatorname{Sq}^2y_4,\
 \operatorname{Sq}^3y_4,\\
 \operatorname{Sq}^4z_5,\
 \operatorname{Sq}^4\operatorname{Sq}^2y_4,\
 \operatorname{Sq}^5\operatorname{Sq}^2y_4.
\end{gathered}
 \tag{D.2.40m}
\]
Its only degree-eleven indecomposable is a positive operation.
On the fiber side the only degree-eleven indecomposables are
$\operatorname{Sq}^5\operatorname{Sq}^2r_4$ and
$\operatorname{Sq}^4\operatorname{Sq}^2s_5$.
The former is an operation of the actual class $r_4$.
For the latter, choose any lift $v_7$ of the permanent fiber
class $\operatorname{Sq}^2s_5$; then
$\operatorname{Sq}^4v_7$ has the required leading filtration term.
The new class $s_5^2$ has degree ten and no degree-one factor
can produce a further degree-eleven indecomposable.
All other filtration terms of degree eleven are products of
positive-degree permanent lifts.  Induction over the Serre
filtration therefore proves the spanning assertion without
requiring every multiplicative or Steenrod extension.
The extra factor $K(\FF_2,5)$ has only the degree-eleven
indecomposables $\operatorname{Sq}^5\operatorname{Sq}^1$ and
$\operatorname{Sq}^4\operatorname{Sq}^2$ of its fundamental
class, so altogether
\[
 H^{11}(P_5U)=
 \sum_{a=1}^{10} H^a(P_5U)H^{11-a}(P_5U)
 +\sum_{j=1}^{11}\operatorname{Sq}^jH^{11-j}(P_5U).
 \tag{D.2.40n}
\]

\medskip\noindent
\textbf{The canonical class detects failure of descent.}
Let $q_4\in H_4U$ be the top homology value of the canonical
Moore map $M(\mathbb Z/2,3)\to U$ from D.2.29.
Its bottom value is $2\alpha$ and has zero mod-two Hurewicz,
whereas two homology suspensions send $q_4$ to $k$.
Thus $q_4$ is primitive and $\beta q_4=0$; all other positive
homology squares vanish by connectivity.
The class $Q^7q_4\in H_{11}U$ is the top $E_4$ operation.
It is primitive, by the diagonal Cartan formula and the
vanishing of its possible mixed terms with units.
The full top Nishida formula, including its Browder correction,
shows that every positive homology square of this class is zero;
see \cite[Chapter~III, Theorems~1.2--1.4,
pp.~216--219]{I_CohenLadaMay1976}.
Explicitly, the potentially nonzero leading coefficients for
indices $1,2,3$ are
$\binom61,\binom52,\binom43$, all even.
Every remaining term involving a positive square of $q_4$,
including the Browder correction, vanishes.  For indices four
and five the remaining operations are unstable; for larger
indices ordinary Steenrod instability on the dual degree gives
zero.  Hence the image of $Q^7q_4$ in $H_{11}(P_5U)$ pairs
trivially with every summand of \textup{(D.2.40n)} and is zero.
On the other hand, homology suspension commutes with the
operations in this range and gives
\[
 \sigma^2(Q^7q_4)=Q^7k=\xi_1k,
 \qquad \langle\mu_{11},Q^7q_4\rangle=1.
 \tag{D.2.40o}
\]
This proves \textup{(D.2.40c)}.
The fiber of $U\to P_5U$ still has mixed Serre terms in the
range at issue.  Failure of descent therefore supplies no
assertion that restriction to that fiber is nonzero, and no
conclusion about the full fourth-suspension attachment follows.
\end{proof}
